\section{引言}\label{sec:intro}

储备池计算只训练一个线性读出层，把输入送进一个固定的非线性动力系统，让它把输入的历史展开成高维特征~\cite{jaeger2001,jaeger2004}。
因为动力系统本身不需要训练，它很适合用物理器件实现，光学是其中研究最多的方向之一~\cite{vandersande2017}。
频率复用的方案把一个光学腔的各根纵模当作储备池的节点，用一个腔、一组波长同时得到几十个特征~\cite{butschek2020,lupo2023}；
光学的下一代储备池计算也沿着类似的思路发展~\cite{cox2024,wang2025ngrc}。
Kerr 微环是这类方案的自然候选：它的四波混频把泵浦的能量转移到几十根梳齿上，提供了强的、片上的非线性，梳齿间隔又宽到可以逐根解复用读出。
近两年已有几种微梳储备池方案：调制泵浦功率、工作在混沌区的光梳~\cite{shishavan2025}，调制泵浦失谐、工作在单孤子态~\cite{cuevas2025}，
调制泵浦相位、工作在腔孤子态~\cite{arabieh2026soliton}，以及用同步脉冲串整体驱动的频率复用腔~\cite{arabieh2026pulse}。
它们有一个共同点：信号只通过一个自由度（泵浦的功率、失谐或相位，或者所有梳齿一起调制）进入器件。
这带来两个问题。在混沌区，梳齿由调制不稳定性从噪声中长出来，同样的输入两次给出的输出可能完全不同（第~\ref{sec:mechanism}~节）；
在孤子区，输出可以重复，但要把器件稳定在孤子存在的区间内，输入能施加的扰动也因此受限。

更根本的问题是，这些工作和软件比较的条件并不统一：工作点怎样选、在哪一次噪声实现上报告、给软件基线设多大的噪声，都各不相同，
而我们发现这几个选择中的任何一个都能让比较的结论翻转。
所以一个基本的问题还没有回答：\emph{在现实噪声下，Kerr 微梳储备池最好能做到什么程度？被什么限制？怎样和软件公平比较？}
只用一个输入自由度的方案离这个上限有多远，也就不清楚。

本文用仿真回答这个问题。为了逼近上限，我们提出谱相位输入掩模：把一个与微环对齐的输入光梳的 16 根梳齿各加上随机的复振幅偏置和调制权重，与泵浦一起送入微环。
输入梳齿本身就落在谐振上，边带由输入相干地“种”下去，所以微环工作在一个输出由输入决定的相干种子区。
我们在同一个仿真器、同样的器件参数和光功率范围、同样的读出、同样的噪声模型、同样的选点预算下，比较了八种编码，
并用多组独立的噪声实现选点和报告。主要结论如下：
\begin{enumerate}
\item \emph{编码决定能否接近上限}（第~\ref{sec:mechanism}、\ref{sec:encodings}~节）。谱相位掩模在 Mackey--Glass 预测上比单泵浦编码和整体调制的光梳好一到两个数量级，
      在信道均衡上小幅领先，在 NARMA-10 上单环不成立。起作用的是梳齿相位的随机性，实现上只需一个调制器加两个谱相位整形器。
\item \emph{上限由什么决定}（第~\ref{sec:noise}、\ref{sec:readout}~节）。理想探测器下的误差下限约一半来自失谐抖动、一半来自散粒噪声，换平台、加功率都降不了它。
      失谐抖动在读出端扣不掉，所以激光的本征线宽是硬指标：本文的器件要求在约 $5~\mathrm{kHz}$ 以下。读出侧的两种改法能把误差再降约 2 倍。
      微环贡献的是非线性特征而不是记忆，记忆要靠外加的延迟线。
\item \emph{上限在软件坐标里的位置}（第~\ref{sec:software}~节）。最好的单环相当于输入接约 7 位模数转换器的 31 节点 ESN；
      在本文考察的任务上，微梳储备池的精度上限低于同规模软件。
\item \emph{设计工具}（第~\ref{sec:gradient}~节）。可微分仿真能在噪声下微调器件参数和掩模，稳定地把起点的误差降低约一半，但结果与扩大后的随机搜索打平，计算量却大得多；
      在本文的参数规模下，单用随机搜索更划算。
\end{enumerate}
第~\ref{sec:methods}~节介绍模型、噪声和评估协议；第~\ref{sec:discussion}~节讨论实验实现、功耗和局限，并简述几个试过但不成立的方向。
